\documentclass[14pt,usepdftitle=false,aspectratio=169]{beamer}
\usepackage{preambule}
\setbeamercolor{structure}{fg=black}
\DeclareMathOperator{\oldvol}{vol}\newcommand{\vol}[1]{\oldvol\l#1\r}
\hypersetup{pdftitle={Théorie algébrique des nombres -- Unités de l’anneau des entiers}}
\title{Théorie algébrique des nombres\\\emph{Unités de l'anneau des entiers}}
\author{}
\date{}
\begin{document}
\begin{frame}
    \titlepage
\end{frame}
\slideq{Famille fondamentale d'unités pour $\calO_K$}{1/4}
\slider{Famille $\l u_1,\cdots,u_{r_1+r_2-1}\r$ de $\calO_K^\times$ qui forme une base de $\calO_K^\times/\mu_K$}{1/4}
\slideq{Théorème des unités de Dirichlet}{2/4}
\slider{Le groupe $\mu_K=\set{x\in\calO_K,\exists n\in\bbN^*,x^n=1}$ est fini et $\calO_K/\mu_K\cong\bbZ^{r_1+r_2-1}$\linebreak$r_1$ le nombre de plongements réels $K\hookrightarrow\bbR$\linebreak$r_2$ le nombre de plongements complexes conjugués $K\hookrightarrow\bbC$}{2/4}
\slideq{$\vol{\operatorname{Log}\l\calO_K^\times\r}$}{3/4}
\slider{$\sqrt{r_1+r_2}R_K$}{3/4}
\slideq{Régulateur de $K$\linebreak$R_K$}{4/4}
\slider{Valeur absolue d'un mineur de taille $r_1+r_2-1$ de la matrice des $\operatorname{Log}{u_i}$ où $\l u_1,\cdots,u_{r_1+r_2-1}\r$ est une famille fondamentale d'unités et $\operatorname{Log}\l u_i\r=\l\sigma_1\l u_i\r,\cdots,\sigma_{r_1+r_2}\l u_i\r\r$ avec $\sigma_j$ les plongements réels et complexes non conjugués}{4/4}
\end{document}